\subsection{Relèvement TPK de la fibre ternaire et partition cylindrique complète}
\label{subsec:elevacion-catalogal-ch}

La fibre de \(729=3^6\) sous-cylindres de raffinement ne provient pas du quotient de transfert \(\mathcal K_{\rm ph}\). Cet opérateur est une représentation ultérieure à mémoire réduite et ne peut engendrer d'historiques. La couverture est construite directement dans le système \((H_k,\rho_k,\operatorname{Ext}_k)\) au moyen des trois prolongements d'un tour de l'holonomie nonadique.

\subsubsection{Inversion de Hensel dans la tour arithmétique et réalisation par histoires}
\label{sec:ch-capacidad-hensel}

Avant de relever les branches en histoires enrichies, il convient de retrouver
l’argument fini qui détermine la capacité de la tour arithmétique. On
incorpore la preuve complète du transducteur non filtré de la matrice ; son
domaine est distingué de l’arbre survivant.

\begin{proposition}[Conjugaison de la tour de Hensel non filtrée]
\label{prop:ch-hensel-bruto}
Soient \(p\) premier, \(d\geq1\) et \(A\in\operatorname{GL}_d(\mathbb Z)\). Pour des vecteurs lignes, en prenant des représentants de résidus dans \(\{0,\ldots,p-1\}^d\), on définit
\[
 x_tA=u_t+px_{t+1},\qquad u_t=x_tA\bmod p.
\]
L'application
\[
 \Psi_T:(\mathbb Z/p^T\mathbb Z)^d\longrightarrow
             (\mathbb F_p^d)^T,\qquad
 x_0\longmapsto(u_0,\ldots,u_{T-1})
\]
est bijective, d’inverse
\[
 x_0\equiv\sum_{t=0}^{T-1}p^tu_tA^{-(t+1)}\pmod{p^T}.
\]
Les bijections commutent avec les réductions et donnent un homéomorphisme
\(\Psi:\mathbb Z_p^d\to(\mathbb F_p^d)^{\mathbb N}\), où
\(\mathbb Z_p:=\varprojlim_T\mathbb Z/p^T\mathbb Z\). L’actualisation
\(x_t\mapsto x_{t+1}\) correspond à supprimer le premier bloc.
\end{proposition}

\begin{proof}
L'inverse entier de \(A\) existe car son déterminant est \(\pm1\). En itérant \(x_t=(u_t+px_{t+1})A^{-1}\), on obtient
\[
 x_0=\sum_{t=0}^{T-1}p^tu_tA^{-(t+1)}+p^Tx_TA^{-T}.
\]
Modulo \(p^T\), un mot détermine une unique classe initiale. Étant donné un
mot arbitraire, fixer \(x_T=0\) et reconstruire à rebours produit une
classe qui l’émet, ce qui prouve la surjectivité. La même formule
prouve la compatibilité lors de la troncature. Le passage à la limite inverse donne
l’homéomorphisme et la relation avec le décalage du ruban.
\end{proof}

Pour \(p=3\) et \(d=6\), l’alphabet de blocs a \(3^6=729\)
éléments. Cette capacité arithmétique n’est pas identifiée par définition à
la réalisabilité de chaque bloc après chaque préfixe survivant. Le
lien avec les arêtes \(\operatorname{Ext}\) est précisément la fonction
du lemme~\ref{lem:elevacion-catalogal-nonadica}. Les deux
opérations que distingue la matrice sont ainsi maintenues : représentation dans la tour brute et
réalisation sous les compatibilités APP--TRIT--TPK.



\subsubsection{Réalisation arête par arête des trois prolongements nonadiques}
\label{subsubsec:tres-vueltas-ext-ch}

Cette sous-section spécifie le modèle autonome engendré par la relation
d’extension TPK qui intervient dans la preuve cardinale. On n’infère pas l’existence
d’une arête par coïncidence de types : on construit ses états initial et final,
sa relation APP, son transport de fibre et la mise à jour unaire du registre.
La construction utilise exclusivement les deux feuillets d’APP, le sélecteur TRIT
et la composition \(\mathsf{Sel}\)--\(\mathsf{Tra}\)--\(\operatorname{Upd}\) du TPK.

\paragraph{Cellules complètes et trois réductions équilibrées.}
Dans la coordinatisation \(D_9=\{1,\ldots,9\}\), on conservera toujours la cellule APP complète
\[
 a_{p,q}:=(p,q;R^\Sigma_{p,q},q^\Sigma_{p,q},
                 R^\Pi_{p,q},q^\Pi_{p,q}),
\]
avec
\[
 p+q=R^\Sigma_{p,q}+9q^\Sigma_{p,q},\qquad
 pq=R^\Pi_{p,q}+9q^\Pi_{p,q}.
\]
Le feuillet actif est enregistré dans une coordonnée séparée ; \(a_{p,q}\) n'est pas scindé en une prétendue cellule additive et une autre multiplicative. Pour \(r\in D_9\), soient
\begin{equation}
 a(r):=
 \begin{cases}
  a_{1,9},&r=1,\\
  a_{1,r-1},&2\le r\le9,
 \end{cases}
 \qquad
 \tau(r):=M_{\rm ph}(r),
 \qquad
 m(r):=\tau(r)\in\{-1,0,+1\}.
 \label{eq:celula-completa-fase-ch}
\end{equation}
Le reste additif de \(a(r)\) fournit ici un représentant APP de la
phase ; cela ne transforme pas \(\Sigma\) en feuillet actif universel. Le sélecteur
dynamique active \(\Sigma\) dans \(r\in\{1,4,7\}\), active \(\Pi\) dans
\(r\in\{2,5,8\}\) et conserve les deux feuillets, sans déplacer le curseur, dans
\(r\in\{3,6,9\}\). Dans les trois cas, on transporte conjointement
\(R^\Sigma,q^\Sigma,R^\Pi,q^\Pi\). On a
\[
 (m(1),\ldots,m(9))=(1,-1,0,1,-1,0,1,-1,0).
\]
Les trois paires ordonnées par la base nonadique sont
\begin{equation}
 \mathfrak p_{+1}:=(1,4),\qquad
 \mathfrak p_{-1}:=(2,5),\qquad
 \mathfrak p_{0}:=(3,6).
 \label{eq:tres-pares-internos-ch}
\end{equation}
Ces paires sont la section normalisée qui prend les deux premières occurrences de chaque type TRIT sous le déplacement \(r\mapsto r+3\) ; le troisième bloc de phases \(7,8,9\) conserve le contrôle de clôture du tour. Leurs deux membres appartiennent à une même fibre de \(M_{\rm ph}\). La réduction équilibrée du chapitre TRIT donne, sans introduire de coefficient cible,
\begin{equation}
 \begin{array}{c|c|c|c}
 \varepsilon&m(\mathfrak p_\varepsilon^{(1)})
                  +m(\mathfrak p_\varepsilon^{(2)})
   &r_3\bigl(2\varepsilon\bigr)&
   \chi(\varepsilon,\varepsilon)\\ \hline
 +1&1+1=(-1)+3\cdot1&-1&+1\\
 0&0+0=0+3\cdot0&0&0\\
 -1&(-1)+(-1)=1+3\cdot(-1)&+1&-1
 \end{array}
 \label{eq:tres-carry-derivados-ch}
\end{equation}
Par conséquent,
\begin{equation}
 \chi(\varepsilon,\varepsilon)
 =\frac{2\varepsilon-r_3(2\varepsilon)}3=\varepsilon.
 \label{eq:carry-no-parametrico-ch}
\end{equation}
Le signe de la branche est ainsi la sortie du cocycle équilibré appliqué à deux
émissions TRIT produites par APP. Chaque paire fournit l’une des trois
sorties de retenue. Par souci de clarté, nous indexerons ensuite la paire par
\(\varepsilon(\mathfrak p):=
\chi(m(\mathfrak p^{(1)}),m(\mathfrak p^{(2)}))\) ; le symbole
\(\varepsilon\) ne fournit pas à \(\operatorname{Upd}^{\rm cal}\) une retenue libre.

L’architecture de phase, le relèvement équilibré et le transport de
mémoire sont réalisés ici au moyen de vingt-sept incidences locales et de leurs
graphes d’extension. La section normalisée fixe des représentants des
fibres de \(M_{\rm ph}\) ; la couverture qui sera démontrée utilise l’existence
de ces représentants et demeure invariante lorsqu’on les remplace par une autre
section compatible.

\paragraph{Données typées des neuf arêtes.}
Soit
\(K_{\rm Carr}:=\langle\kappa_{\rm Carr}\rangle_{\rm Ab}\) le module de
retenue. Les actions sont typées par
\[
 H_{\rm Carr}:\mathscr P_{\rm TPK}^{\rm op}
   \longrightarrow\operatorname{End}(K_{\rm Carr}),\qquad
 Y_{\rm Carr}:\mathscr P_{\rm TPK}
   \longrightarrow\operatorname{End}(K_{\rm Carr}).
\]
Dans ce modèle engendré, on prend les actions triviales \(H_{\rm Carr}(e)=Y_{\rm Carr}(e)=\operatorname{id}_{K_{\rm Carr}}\) : le feuillet actif et le chemin sont conservés comme coordonnées propres du registre et n'agissent pas sur le terme libre de retenue. Soit en outre \((\mathsf M_d,\jmath_{d',d})\) le système inductif gradué de mémoire. Pour ne pas confondre deux représentants identifiés par la colimite, l'état conserve en outre le degré comme coordonnée et l'on utilise l'espace porteur étiqueté
\[
 \widetilde{\mathsf M}:=
 \bigsqcup_{d\ge1}\bigl(\{d\}\times\mathsf M_d\bigr),
 \qquad
 \operatorname{gr}_{\mathsf M}(d,\mu)=d.
\]
La représentation dans la mémoire limite est donnée par
\[
\begin{aligned}
 \mathsf M&:=\varinjlim_d\mathsf M_d,&
 \jmath_d&:\mathsf M_d\longrightarrow\mathsf M,\\
 u&:=\jmath_d(u_d),&
 \iota_{\mathsf M}:\mathbb Z&\longrightarrow\mathsf M,
 \qquad n\longmapsto nu .
\end{aligned}
\]
Ici, \(\jmath_d\) est l'application canonique de la colimite et les \(u_d\) forment une collection compatible. Sa composante libre permet le lecteur \(\deg_u(nu)=n\). Les projections de configuration, phase, mémoire, retenue, route et survie seront typées sur le support calibré construit récursivement plus bas ; aucun espace ambiant ne sera importé pour elles. Pour chaque arête construite \(e\), la composante calibrée \(\operatorname{Upd}^{\rm cal}_e:=\mathsf{Mem}^{\rm cal}_e\) désignera la troisième application de la factorisation ; sa loi est fixée champ par champ ci-dessous et n'est pas présupposée comme restriction d'une relation ambiante.

Fixons la profondeur initiale \(k_0=1\), une configuration complète
\(q_\star\in\mathcal Q_{\rm obs}\) de phase un et
\[
 q_{n,r}:=F_{\rm obs}^{9n+r-1}(q_\star),\qquad
 q_{n,10}:=q_{n+1,1}.
\]
Comme le calendrier observable est périodique, sa valeur \(q_{n,r}\) ne
distingue pas à elle seule des apparitions séparées par douze tours. Le modèle
calibré emploie donc le relèvement étiqueté
\[
 \widetilde q_{n,r}:=(q_{n,r},n,r),\qquad
 \widetilde q_{n,10}:=\widetilde q_{n+1,1},\qquad
 \pi_{\rm obs}(\widetilde q_{n,r})=q_{n,r}.
\]
La dynamique élevée \(\widetilde F_{\rm obs}(\widetilde q_{n,r})
=\widetilde q_{n,r+1}\) se projette sur \(F_{\rm obs}\) ; l'étiquette ne modifie pas la phase physique, mais distingue ses niveaux d'apparition. On étiquette la copie de la cellule à chaque niveau au moyen de
\[
 a_{n,r}:=(n,r;a(r)),\qquad
 (n,r)^+:=
 \begin{cases}
  (n,r+1),&1\le r<9,\\
  (n+1,1),&r=9.
 \end{cases}
\]
L’étiquette évite seulement d’identifier deux apparitions d’une même cellule ; les
deux divisions euclidiennes et leurs quatre coordonnées sont celles de \(a(r)\).
Pour le tour \(n\ge0\) et la phase \(r\in D_9\), le symbole
\[
 e_{\varepsilon,n,r}:\widetilde q_{n,r}
 \longrightarrow\widetilde q_{n,r+1},
 \qquad r^+:=1+(r\bmod9),
\]
occupe le niveau \(d_{n,r}:=k_n+r-1\), où \(k_n:=k_0+9n\). Son registre local contient la cellule complète \(a(r)\), le feuillet actif déterminé par \(M_{\rm ph}(r)\), le type \(\tau(r)\), la marque d'appartenance à \(\mathfrak p_\varepsilon\) et, lorsqu'il atteint le second membre de cette paire, l'identité certifiée \eqref{eq:tres-carry-derivados-ch}. Nous définissons ses opérateurs uniquement à partir de ce registre :
\begin{align}
 \mathsf C_{\mathsf M}(e_{\varepsilon,n,r})
   &=\jmath_{d_{n,r}+1,d_{n,r}},&
 \kappa_{\mathsf M}(e_{\varepsilon,n,r})
   &=\mathbf1_{\{r=9\}}u_{d_{n,r}+1},
 \label{eq:memoria-arista-ch}\\
       \operatorname{Carr}(e_{\varepsilon,n,r})
   &=d_{\varepsilon,r}\kappa_{\rm Carr},&
 d_{\varepsilon,r}
   &:=\mathbf1_{\{r=\mathfrak p_\varepsilon^{(2)}\}}
       \chi\!\left(m(\mathfrak p_\varepsilon^{(1)}),
                    m(\mathfrak p_\varepsilon^{(2)})\right).
 \label{eq:carry-arista-derivado-ch}
\end{align}
Pour les identités et pour toute composition admissible, on définit en outre
\begin{align}
 \mathsf C_{\mathsf M}(\operatorname{id}_{\widetilde q};d)
   &:=\operatorname{id}_{\mathsf M_d},&
 \kappa_{\mathsf M}(\operatorname{id}_{\widetilde q};d)&:=0,
 \label{eq:memoria-identidad-ch}\\
 \mathsf C_{\mathsf M}(e_t\cdots e_1)
   &:=\mathsf C_{\mathsf M}(e_t)\circ\cdots\circ
      \mathsf C_{\mathsf M}(e_1),&
 \kappa_{\mathsf M}(e_2e_1)
   &:=\mathsf C_{\mathsf M}(e_2)
        \bigl(\kappa_{\mathsf M}(e_1)\bigr)
      +\kappa_{\mathsf M}(e_2).
 \label{eq:cociclo-memoria-calibrado-ch}
\end{align}
L’identité porte le degré \(d\) car une même configuration observable
peut réapparaître à des profondeurs distinctes. La dernière égalité, itérée, définit
\(\kappa_{\mathsf M}(e_t\cdots e_1)\) pour toute longueur et est la loi de
cocycle du système direct ; les inclusions satisfont
\(\jmath_{d+2,d+1}\jmath_{d+1,d}=\jmath_{d+2,d}\).
On fixe en outre
\(s(e_{\varepsilon,n,r})=\top\). Les vingt-sept entrées suivantes
certifient les coordonnées qui distinguent les trois prolongements ; les champs
universels du transport sont spécifiés immédiatement après. Les symboles \(I\) et \(II\)
désignent respectivement la première émission conservée et la seconde émission
qui effectue la réduction ; un tiret indique un accroissement de retenue nul, mais non
une absence d’arête.
\begin{center}
\small
\setlength{\tabcolsep}{3.2pt}
\begin{tabular}{c c cc cc cc c}
\toprule
\(r\)&\(M_{\rm ph}(r)\)&\multicolumn{2}{c}{\(\mathfrak p_{+1}\)}&
\multicolumn{2}{c}{\(\mathfrak p_0\)}&
\multicolumn{2}{c}{\(\mathfrak p_{-1}\)}&
\(\Delta\operatorname{mem}\)\\
& &marque&\(d_{+1,r}\)&marque&\(d_{0,r}\)&marque&\(d_{-1,r}\)&\\
\midrule
1&\(+1\)&I&0&--&0&--&0&0\\
2&\(-1\)&--&0&--&0&I&0&0\\
3&0&--&0&I&0&--&0&0\\
4&\(+1\)&II&\(+1\)&--&0&--&0&0\\
5&\(-1\)&--&0&--&0&II&\(-1\)&0\\
6&0&--&0&II&0&--&0&0\\
7&\(+1\)&--&0&--&0&--&0&0\\
8&\(-1\)&--&0&--&0&--&0&0\\
9&0&--&0&--&0&--&0&\(u\)\\
\bottomrule
\end{tabular}
\end{center}
Ainsi, \(d_{\varepsilon,r}\) est calculé à partir de deux sorties du sélecteur et
est distinct des quotients \(q^\Sigma,q^\Pi\) de la cellule. Le registre
conserve également la paire marquée et le résidu équilibré ; en particulier, la
branche nulle n’est pas confondue avec une absence d’événement.

La loi de composition de la retenue employée par la catégorie TPK est
\begin{equation}
 \operatorname{Carr}(\operatorname{id}_{\widetilde q})=0,\qquad
 \operatorname{Carr}(e_2e_1)
 =H_{\rm Carr}(e_1)\bigl(\operatorname{Carr}(e_2)\bigr)
  +Y_{\rm Carr}(e_2)\bigl(\operatorname{Carr}(e_1)\bigr).
 \label{eq:ley-carry-correcta-ch}
\end{equation}
Sur les identités et les chemins, ils se prolongent fonctoriellement :
\[
\begin{aligned}
 H_{\rm Carr}(\operatorname{id})
   &=Y_{\rm Carr}(\operatorname{id})=\operatorname{id},\\
 H_{\rm Carr}(e_t\cdots e_1)
   &=H_{\rm Carr}(e_1)\circ\cdots\circ H_{\rm Carr}(e_t),\\
 Y_{\rm Carr}(e_t\cdots e_1)
   &=Y_{\rm Carr}(e_t)\circ\cdots\circ Y_{\rm Carr}(e_1).
\end{aligned}
\]
Dans ce modèle, toutes ces actions restent l'identité, mais leur typage explicite quel endomorphisme agit sur chaque terme de \eqref{eq:ley-carry-correcta-ch}.
Sur les arêtes calibrées, \(H_{\rm Carr}=Y_{\rm Carr}=\operatorname{id}\), si bien que la composition additionne les accroissements déjà produits. L'équation \eqref{eq:ley-carry-correcta-ch}, et non une identification entre orientation et quotient, gouverne le registre accumulé.

\paragraph{Construction indépendante des tuples de transition.}
Un état de la fibre sera écrit, dans l’ordre de ses champs pertinents,
\[
 x=(a;q;\tau,o,h;b_{\rm ref};
       R^\Sigma,q^\Sigma,R^\Pi,q^\Pi;c;
       \operatorname{Sig};C,\partial C;p_{\APP},p_{\TRIT};
       \lambda;\operatorname{gr}_{\mathsf M};\operatorname{mem};
       s;\operatorname{Led}).
\]
La notation de mise à jour \(x[f\leftarrow v]\) ne change que le champ
indiqué. Pour \(h\in\{(\Sigma,\Pi),(\Pi,\Sigma)\}\), définissons
\[
 h_r(h)=
 \begin{cases}
  (\Sigma,\Pi),&r\in\{1,4,7\},\\
  (\Pi,\Sigma),&r\in\{2,5,8\},\\
 h,&r\in\{3,6,9\}.
 \end{cases}
\]
La coordonnée complète de feuillet et d’arithmétique est abrégée par
\[
 \widehat h(x):=
 \bigl(h(x),R^\Sigma(x),q^\Sigma(x),R^\Pi(x),q^\Pi(x)\bigr).
\]
Soit \(\widehat{\mathcal H}^{\rm cal}_{\widetilde q_{n,r}}\) l’ensemble des
quintuplets
\[
 \bigl(h,R^\Sigma_{a_{n,r}},q^\Sigma_{a_{n,r}},
          R^\Pi_{a_{n,r}},q^\Pi_{a_{n,r}}\bigr),
 \qquad
 h\in\{(\Sigma,\Pi),(\Pi,\Sigma)\},
\]
dont les quatre données arithmétiques sont les deux divisions euclidiennes de la cellule
APP étiquetée \(a_{n,r}\). Cette définition directe ne présuppose pas d’état
sur un porteur qui n’est pas encore construit. Soient
\(\mathsf{Hist}^{\rm cal}_{\widetilde q}\) l’ensemble des chemins finis
admissibles du graphe d’arêtes étiquetées qui aboutissent à
\(\widetilde q\), y compris l’identité, et
\(\mathsf B^{\rm cal}_{\widetilde q}\) l’ensemble de leurs suites
ordonnées d’extrémités. Nous définissons
\[
 \mathsf S^{\rm cal}_{\widetilde q}:=
 \widehat{\mathcal H}^{\rm cal}_{\widetilde q}
 \times\{o_{\rightarrow},o_{\circ},o_{\leftarrow}\}
 \times K_{\rm Carr}\times
 \mathsf{Hist}^{\rm cal}_{\widetilde q}\times
 \mathsf B^{\rm cal}_{\widetilde q}\times\mathsf M .
\]
La signature calibrée n’est pas postulée : c’est l’application de réunion typée
\[
 \operatorname{Sig}^{\rm cal}_{\widetilde q}:
 \widehat{\mathcal H}^{\rm cal}_{\widetilde q}
 \times\{o_{\rightarrow},o_{\circ},o_{\leftarrow}\}
 \times K_{\rm Carr}\times
 \mathsf{Hist}^{\rm cal}_{\widetilde q}\times
 \mathsf B^{\rm cal}_{\widetilde q}\times\mathsf M
 \longrightarrow\mathsf S^{\rm cal}_{\widetilde q},
\quad
 (\widehat h,o,c,\lambda,\partial C,\mu)
 \longmapsto(\widehat h,o,c,\lambda,\partial C,\mu).
\]
De même, \(\pi_{\mathcal F}\) désigne la projection qui élimine uniquement les deux premières coordonnées \((a;q)\) du tuple d'état ; elle conserve \(\tau,o,h,b_{\rm ref}\), les quatre données euclidiennes, la retenue, la signature, le cylindre, la frontière, les préfixes, le chemin, le degré de mémoire, la mémoire, la survie et le registre.
L'orientation correspondante est fixée par \(o(+1)=o_{\rightarrow}\), \(o(0)=o_{\circ}\) et \(o(-1)=o_{\leftarrow}\). La fibre APP calibrée sur \(\widetilde q_{n,r}\) contiendra la copie étiquetée \(a_{n,r}\) déjà définie. L'état initial est fixé sans coordonnées libres :
\[
 \widehat h_\star:=
 \bigl((\Sigma,\Pi),R^\Sigma_{1,9},q^\Sigma_{1,9},
       R^\Pi_{1,9},q^\Pi_{1,9}\bigr).
\]
\begin{equation}
\begin{aligned}
 x_\star:=\bigl(&a_{0,1};\widetilde q_{0,1};
 +1,o_{\rightarrow},(\Sigma,\Pi);\bot;
 R^\Sigma_{1,9},q^\Sigma_{1,9},R^\Pi_{1,9},q^\Pi_{1,9};0;\\
 &\operatorname{Sig}^{\rm cal}_{\widetilde q_{0,1}}
       (\widehat h_\star,o_{\rightarrow},0,
        \operatorname{id}_{\widetilde q_{0,1}},\varnothing,
        \jmath_{k_0}(0_{\mathsf M_{k_0}}));
 \varnothing,\varnothing;
 (a_{0,1}),(+1);\operatorname{id}_{\widetilde q_{0,1}};
 k_0;0_{\mathsf M_{k_0}};\top;\varnothing\bigr).
\end{aligned}
\label{eq:estado-inicial-calibrado-ch}
\end{equation}
Le tuple ne contient pas de coordonnées indéterminées et satisfait
\(s(x_\star)=\top\). Soit \(G_0:=\{x_\star\}\).
Supposant \(G_n\) construit, fixons \(x\in G_n\),
\(\varepsilon\in\{-1,0,+1\}\) et
\(x_{\varepsilon,0}(x):=x\). Pour \(r=1,\ldots,9\), on construit d’abord
le registre élémentaire
\begin{equation}
\begin{aligned}
 \ell_{\varepsilon,n,r}:=\bigl(&e_{\varepsilon,n,r},a_{n,r},a_{(n,r)^+},
 \widetilde q_{n,r},\widetilde q_{n,r+1},M_{\rm ph}(r),\\
 &h_r(h(x_{\varepsilon,r-1}(x))),o(M_{\rm ph}(r)),\mathfrak p_\varepsilon,
 d_{\varepsilon,r},\\
 &\mathsf C_{\mathsf M}(e_{\varepsilon,n,r}),
 \kappa_{\mathsf M}(e_{\varepsilon,n,r}),
 \operatorname{Carr}(e_{\varepsilon,n,r})\bigr),
\end{aligned}
 \label{eq:registro-elemental-ext-ch}
\end{equation}
et on pose
\(\operatorname{Led}^{\rm new}_{\varepsilon,n,r}:=
\operatorname{Led}(x_{\varepsilon,r-1}(x))^\frown
\ell_{\varepsilon,n,r}\). On construit alors le tuple
\begin{equation}
 y_{\varepsilon,r}(x):=
 x_{\varepsilon,r-1}(x)
 [\tau\leftarrow M_{\rm ph}(r),\
  h\leftarrow h_r(h(x_{\varepsilon,r-1}(x))),\
  o\leftarrow o(M_{\rm ph}(r)),\
  b_{\rm ref}\leftarrow\mathfrak p_\varepsilon],
 \label{eq:tupla-sel-cal-ch}
\end{equation}
où la sélection laisse intacts la cellule, les quotients, la retenue, la route, la mémoire,
le cylindre, la frontière, la survie et le registre ; la coordonnée
\(b_{\rm ref}\) explicite la réduction équilibrée choisie. On ne présuppose pas ici
l’appartenance à une relation ambiante.

On définit ensuite \(z_{\varepsilon,r}(x)\) champ par champ : sa
configuration est \(\widetilde q_{n,r+1}\) ; sa cellule complète est
\(a_{(n,r)^+}\), avec les
quatre résidus et quotients recalculés par les deux divisions euclidiennes ;
il conserve le feuillet sélectionné et l’orientation ; et il effectue les six
mises à jour suivantes, dans lesquelles seul l’argument commun \(x\) est omis :
\begin{equation}
\begin{aligned}
 C(z_{\varepsilon,r})&:=C(y_{\varepsilon,r})^\frown
       e_{\varepsilon,n,r},\\
 \partial C(z_{\varepsilon,r})&:=\partial C(y_{\varepsilon,r})^\frown
       (\widetilde q_{n,r},\widetilde q_{n,r+1}),\\
 p_{\APP}(z_{\varepsilon,r})&:=p_{\APP}(y_{\varepsilon,r})^\frown
       a_{(n,r)^+},\\
 p_{\TRIT}(z_{\varepsilon,r})&:=p_{\TRIT}(y_{\varepsilon,r})^\frown
       \tau(r),\\
 \operatorname{Led}(z_{\varepsilon,r})&:=
       \operatorname{Led}^{\rm new}_{\varepsilon,n,r},\\
 \lambda(z_{\varepsilon,r})&:=e_{\varepsilon,n,r}
       \lambda(y_{\varepsilon,r}),
\end{aligned}
 \label{eq:tupla-tra-cal-ch}
\end{equation}
et maintient encore les champs \(c,\operatorname{gr}_{\mathsf M},\operatorname{mem},s\) de \(y_{\varepsilon,r}(x)\). Sa signature, en revanche, est à nouveau typée dans la configuration de destination au moyen de la présignature transportée
\begin{equation}
\begin{aligned}
 \operatorname{Sig}(z_{\varepsilon,r}(x))
  :=\operatorname{Sig}^{\rm cal}_{\widetilde q_{n,r+1}}
       \Bigl(&\widehat h(z_{\varepsilon,r}(x)),
             o(z_{\varepsilon,r}(x)),c(y_{\varepsilon,r}(x)),
             \lambda(z_{\varepsilon,r}(x)),\\[-2pt]
             &\partial C(z_{\varepsilon,r}(x)),
             \jmath_{d_{n,r}}
                (\operatorname{mem}_{d_{n,r}}
                   (y_{\varepsilon,r}(x)))\Bigr).
\end{aligned}
\label{eq:prefirma-transporte-cal-ch}
\end{equation}
Ainsi, et seulement dans
\(\mathsf{Tra}^{\rm cal}\), on ajoute l’arête, ses extrémités, la marque
\((\mathfrak p_\varepsilon,r,d_{\varepsilon,r})\) et ses cocycles au registre des
incidences.

L'état initial satisfait
\[
 c(x_\star)=0=\operatorname{Carr}
   (\operatorname{id}_{\widetilde q_{0,1}})
   =\operatorname{Carr}(\lambda(x_\star)).
\]
Dans les états complets \(x_{\varepsilon,r-1}(x)\), et aussi après la
sélection \(y_{\varepsilon,r}(x)\), on conserve l’invariant typé
\(c(v)=\operatorname{Carr}(\lambda(v))\), puisque
\(\mathsf{Sel}^{\rm cal}\) ne modifie aucun de ces deux champs. Le
transport fait avancer la route mais laisse la retenue en attente ; le tuple de
préactualisation satisfait donc exactement
\[
 c(z_{\varepsilon,r}(x))=c(y_{\varepsilon,r}(x))
 =\operatorname{Carr}(\lambda(y_{\varepsilon,r}(x))).
\]
L'application \(\operatorname{Upd}^{\rm cal}\) restaure alors \(c(x_{\varepsilon,r}(x))=
\operatorname{Carr}(\lambda(x_{\varepsilon,r}(x)))\) au moyen de la loi de composition, sans attribuer cet invariant à l'état intermédiaire \(z\).

Enfin, on applique la fonction unaire
\(\operatorname{Upd}^{\rm cal}_{e_{\varepsilon,n,r}}
=\mathsf{Mem}^{\rm cal}_{e_{\varepsilon,n,r}}\) à ce tuple. Son
résultat \(x_{\varepsilon,r}(x)\) est donné par
\begin{equation}
\begin{aligned}
 c\bigl(x_{\varepsilon,r}(x)\bigr)
   &=\operatorname{Carr}\bigl(\lambda(z_{\varepsilon,r}(x))\bigr)\\
   &=H_{\rm Carr}(\lambda(y_{\varepsilon,r}(x)))
       \bigl(\operatorname{Carr}(e_{\varepsilon,n,r})\bigr)
     +Y_{\rm Carr}(e_{\varepsilon,n,r})
       \bigl(c(y_{\varepsilon,r}(x))\bigr),\\
 \lambda\bigl(x_{\varepsilon,r}(x)\bigr)
   &=\lambda(z_{\varepsilon,r}(x)),\\
 \operatorname{gr}_{\mathsf M}
      \bigl(x_{\varepsilon,r}(x)\bigr)
   &=d_{n,r}+1,\\
 \operatorname{mem}_{d_{n,r}+1}
       \bigl(x_{\varepsilon,r}(x)\bigr)
   &=\mathsf C_{\mathsf M}(e_{\varepsilon,n,r})
       \operatorname{mem}_{d_{n,r}}(z_{\varepsilon,r}(x))
     +\kappa_{\mathsf M}(e_{\varepsilon,n,r}),\\
 s\bigl(x_{\varepsilon,r}(x)\bigr)
   &=s(e_{\varepsilon,n,r})\wedge s(z_{\varepsilon,r}(x)),\\
 \operatorname{Sig}\bigl(x_{\varepsilon,r}(x)\bigr)
   &=\operatorname{Sig}^{\rm cal}_{\widetilde q_{n,r+1}}
       \Bigl(\widehat h(z_{\varepsilon,r}(x)),
             o(z_{\varepsilon,r}(x)),
             c(x_{\varepsilon,r}(x)),
             \lambda(z_{\varepsilon,r}(x)),
             \partial C(z_{\varepsilon,r}(x)),\\[-2pt]
   &\hspace{112pt}
             \jmath_{d_{n,r}+1}
              \bigl(\operatorname{mem}_{d_{n,r}+1}
                 (x_{\varepsilon,r}(x))\bigr)\Bigr).
\end{aligned}
\label{eq:upd-unaria-tres-vueltas-ch}
\end{equation}
Tous les champs non énumérés dans \eqref{eq:upd-unaria-tres-vueltas-ch} coïncident avec ceux de \(z_{\varepsilon,r}(x)\). En particulier, la signature est calculée avec le feuillet, l'orientation, la frontière et le chemin déjà fixés par \(\mathsf{Sel}^{\rm cal}\) et \(\mathsf{Tra}^{\rm cal}\), et avec la retenue et la mémoire calculées dans les lignes précédentes ; elle n'est pas définie circulairement à partir de l'état final. Ni \(\varepsilon\), ni \(d_{\varepsilon,r}\), ni \(u\) ne sont des arguments supplémentaires de \(\operatorname{Upd}^{\rm cal}\) : la fonction lit les trois données de l'arête déjà construite.

\paragraph{Récurrence simultanée des niveaux.}
Une fois les neuf étapes précédentes achevées pour un \(G_n\) déjà construit,
on définit immédiatement
\begin{equation}
 \gamma_{\varepsilon,k_n}(x):=x_{\varepsilon,9}(x),\qquad
 G_{n+1}:=\{\gamma_{\varepsilon,k_n}(x):
       x\in G_n,\ \varepsilon\in\TRIT\}.
 \label{eq:injerto-total-tres-vueltas-ch}
\end{equation}
\label{eq:tres-vueltas-tpk-ch}
La base \(G_0=\{x_\star\}\), les tuples intermédiaires et
\eqref{eq:injerto-total-tres-vueltas-ch} constituent une unique définition par
récurrence sur \(n\). Par conséquent, aucun porteur ultérieur n’est utilisé pour
supposer l’existence de \(G_n\) : le niveau suivant n’est introduit
qu’après avoir construit toutes ses arêtes à partir du niveau précédent.

\paragraph{Relation d’extension TPK et appartenance effective.}
Soit \(\mathcal E^{\rm cal}_{\widetilde q}\) la plus petite collection étiquetée contenant \(x_\star\) et, à chaque étape de la récurrence précédente, les quatre tuples \(x_{\varepsilon,r-1}(x),\allowbreak y_{\varepsilon,r}(x),\allowbreak
z_{\varepsilon,r}(x),\allowbreak x_{\varepsilon,r}(x)\) dans leurs configurations de départ et d'arrivée. Soit également \(\mathcal A^{\rm cal}_{\widetilde q_{n,r}}:=\{a_{n,r}\}\), et posons
\[
 \widetilde{\mathcal Q}_{\rm obs}:=
   \{\widetilde q_{n,r}:n\ge0,\ 1\le r\le9\},
 \qquad
 \mathcal E^{\rm cal}:=\bigsqcup_{\widetilde q}
     \mathcal E^{\rm cal}_{\widetilde q},
 \qquad
 \mathcal F^{\rm cal}_{\widetilde q}:=
     \pi_{\mathcal F}(\mathcal E^{\rm cal}_{\widetilde q}).
\]
Pour chaque profondeur \(d\), soit
\(\mathcal E^{\rm cal}_d:=
\operatorname{gr}_{\mathsf M}^{-1}(d)\subset\mathcal E^{\rm cal}\). La
coordonnée explicite de degré, et non la seule appartenance d’un représentant à
une image du système direct, rend ces porteurs disjoints. Sur eux,
les projections sont typées
\[
\begin{aligned}
 \operatorname{cfg}&:\mathcal E^{\rm cal}\longrightarrow
     \widetilde{\mathcal Q}_{\rm obs},\\
 g(x)&:=\operatorname{fase}\!
       \bigl(\pi_{\rm obs}(\operatorname{cfg}(x))\bigr)\in C_9,\\
 c&:\mathcal E^{\rm cal}\longrightarrow K_{\rm Carr},\\
 s&:\mathcal E^{\rm cal}\longrightarrow\{\bot,\top\},\\
 \lambda&:\mathcal E^{\rm cal}_{\widetilde q}
       \longrightarrow\mathsf{Hist}^{\rm cal}_{\widetilde q},\\
 \operatorname{gr}_{\mathsf M}&:\mathcal E^{\rm cal}
       \longrightarrow\mathbb N_{\ge1},\\
 \operatorname{mem}_d&:\mathcal E^{\rm cal}_d
       \longrightarrow\mathsf M_d,\\
 \operatorname{mem}(x)&:=
       \jmath_d\bigl(\operatorname{mem}_d(x)\bigr)\in\mathsf M
       \quad(x\text{ de profondeur }d).
\end{aligned}
\]
Ici, \(\operatorname{cfg}(x)\) est la seconde coordonnée du tuple et
\(\operatorname{fase}\) lit la phase du calendrier observable. Pour les trois
classes de porteur, on fixe les diagonales
\[
\begin{aligned}
 \Delta_{\APP,\widetilde q}
   &:=\{(a,a):a\in\mathcal A^{\rm cal}_{\widetilde q}\},\\
 \Delta_{{\rm fib},\widetilde q}
   &:=\{(f,f):f\in\mathcal F^{\rm cal}_{\widetilde q}\},\\
 \Delta_{{\rm enr},\widetilde q}
   &:=\{(x,x):x\in\mathcal E^{\rm cal}_{\widetilde q}\}.
\end{aligned}
\]
En particulier,
\[
 a_{0,1}\in\mathcal A^{\rm cal}_{\widetilde q_{0,1}},
 \qquad
 x_\star\in\mathcal E^{\rm cal}_{\widetilde q_{0,1}}.
\]
Sur ces espaces porteurs, on définit par leurs graphes complets, et non au moyen d'implications nécessaires, les relations
\begin{align}
 \mathsf{Sel}^{\rm cal}_{e_{\varepsilon,n,r}}
   &:=\{(x_{\varepsilon,r-1}(x),y_{\varepsilon,r}(x)):
          x\in G_n\},\\
 \mathsf{Tra}^{\rm cal}_{e_{\varepsilon,n,r}}
   &:=\{(y_{\varepsilon,r}(x),z_{\varepsilon,r}(x)):
          x\in G_n\},\\
 \operatorname{Ext}^{\APP,\rm cal}_{e_{\varepsilon,n,r}}
   &:=\{(a_{n,r},a_{(n,r)^+})\},
 \label{eq:ext-app-calibrada-ch}\\
 \operatorname{Ext}^{\rm fib,cal}_{e_{\varepsilon,n,r}}
   &:=\{(\pi_{\mathcal F}x_{\varepsilon,r-1}(x),
          \pi_{\mathcal F}x_{\varepsilon,r}(x)):x\in G_n\}.
 \label{eq:ext-fibra-calibrada-ch}
\end{align}
L'extension enrichie synchronisée est
\begin{equation}
 \operatorname{Ext}^{\rm enr,cal}_{e_{\varepsilon,n,r}}
 :=\{(x_{\varepsilon,r-1}(x),x_{\varepsilon,r}(x)):
       x\in G_n\}.
 \label{eq:ext-enriquecida-calibrada-ch}
\end{equation}
La troisième application est définie par
\begin{equation}
 \operatorname{graph}
 \bigl(\operatorname{Upd}^{\rm cal}_{e_{\varepsilon,n,r}}\bigr)
 :=\{(z_{\varepsilon,r}(x),x_{\varepsilon,r}(x)):
       x\in G_n\},
 \label{eq:grafo-upd-calibrado-ch}
\end{equation}
où la seconde composante est exactement celle de
\eqref{eq:upd-unaria-tres-vueltas-ch}. Par conséquent,
\begin{equation}
 \mathcal U^{(1),\rm cal}_{e_{\varepsilon,n,r}}
 :=\operatorname{Upd}^{\rm cal}_{e_{\varepsilon,n,r}}\circ
   \mathsf{Tra}^{\rm cal}_{e_{\varepsilon,n,r}}\circ
   \mathsf{Sel}^{\rm cal}_{e_{\varepsilon,n,r}}
 =\{(x_{\varepsilon,r-1}(x),x_{\varepsilon,r}(x)):
      x\in G_n\}.
 \label{eq:transicion-cal-en-tpk-ch}
\end{equation}
Chaque opérateur est indexé par l’arête
\(e_{\varepsilon,n,r}\) produite par la réduction
\(\mathfrak p_\varepsilon\), et non par un trit fourni de l’extérieur. La
coordonnée \(b_{\rm ref}\) distingue en outre leurs codomaines. Par conséquent,
\(\mathsf{Sel}^{\rm cal}_{e}\),
\(\mathsf{Tra}^{\rm cal}\) et \(\operatorname{Upd}^{\rm cal}\) sont des fonctions
sur leurs domaines respectifs, tandis que
l’union
\(\bigcup_{\varepsilon\in\TRIT}
\mathcal U^{(1),\rm cal}_{e_{\varepsilon,n,r}}\) conserve délibérément
les trois sorties de la trisection.

Pour \(\bullet\in\{\APP,{\rm fib},{\rm enr}\}\), les identités et les compositions ne restent pas implicites. On définit
\begin{align}
 \operatorname{Ext}^{\bullet,\rm cal}_{\operatorname{id}_{\widetilde q}}
   &:=\Delta_{\bullet,\widetilde q},\\
 \operatorname{Ext}^{\bullet,\rm cal}_{e_t\cdots e_1}
   &:=\operatorname{Ext}^{\bullet,\rm cal}_{e_t}\circ\cdots\circ
      \operatorname{Ext}^{\bullet,\rm cal}_{e_1}
 \label{eq:ext-palabra-calibrada-ch}
\end{align}
pour tout mot admissible d’arêtes. Les deux clôtures de chemins APP et de
fibre sont formées séparément,
\begin{equation}
 \operatorname{Path}(\operatorname{Ext}^{\APP,\rm cal}),\qquad
 \operatorname{Path}(\operatorname{Ext}^{\rm fib,cal}),
 \label{eq:dos-clausuras-ext-ch}
\end{equation}
et \(\mathsf{Tra}^{\rm cal}\) les synchronise arête par arête.

Enfin, les niveaux d'états complets et leurs troncatures sont fixés par la récurrence elle-même :
\begin{align}
 H^{\rm cal}_{d_{n,r}}
   &:=\{x_{\varepsilon,r-1}(x):
          x\in G_n,\ \varepsilon\in\TRIT\},\\
 H^{\rm cal}_{d_{n,r}+1}
   &:=\{x_{\varepsilon,r}(x):
          x\in G_n,\ \varepsilon\in\TRIT\},\\
 \rho^{\rm cal}_{d_{n,r}+1,d_{n,r}}
   \bigl(x_{\varepsilon,r}(x)\bigr)
   &:=x_{\varepsilon,r-1}(x).
 \label{eq:truncamiento-calibrado-definido-ch}
\end{align}
Le dernier registre conserve \((n,r,\mathfrak p_\varepsilon)\), de sorte que
le prédécesseur du membre de gauche est unique et \(\rho^{\rm cal}\) est bien
définie. Pour \(b>a\), on pose
\(\rho^{\rm cal}_{b,a}:=\rho^{\rm cal}_{a+1,a}\circ\cdots\circ
\rho^{\rm cal}_{b,b-1}\). Cette application retrouve l’état précédent
complet, y compris sa signature ; elle ne prétend pas inverser une coordonnée écrasée
sans consulter le registre.

\begin{proposition}[Réalisation autonome du TPK au moyen de transitions explicites]
\label{prop:modelo-tpk-calibrado-ch}
La collection
\[
\begin{aligned}
 \mathfrak T_{\rm TPK}^{\rm cal}:=\bigl(
 &(\widetilde q_{n,r})_{n,r},\pi_{\rm obs},
   \widetilde F_{\rm obs},M_{\rm ph},\\
 &(\mathcal A_{\widetilde q}^{\rm cal})_{\widetilde q},
   (\mathcal E_{\widetilde q}^{\rm cal})_{\widetilde q},
   (\mathcal F_{\widetilde q}^{\rm cal})_{\widetilde q},
   \operatorname{Ext}^{\APP,\rm cal},
   \operatorname{Ext}^{\rm fib,cal},
   \operatorname{Ext}^{\rm enr,cal},\rho^{\rm cal},\\
 &\mathsf{Sel}^{\rm cal},\mathsf{Tra}^{\rm cal},
   \operatorname{Upd}^{\rm cal},\mathsf C_{\mathsf M},
   \kappa_{\mathsf M},\operatorname{gr}_{\mathsf M},
   H_{\rm Carr},Y_{\rm Carr},\\
 &\operatorname{Sig}^{\rm cal},\pi_{\mathcal F}\bigr).
\end{aligned}
\]
est un modèle enrichi du TPK. En particulier, chaque paire du membre de droite de \eqref{eq:transicion-cal-en-tpk-ch} est une arête TPK effective, et non la conséquence inférée d'une condition seulement nécessaire.
\end{proposition}

\begin{proof}
Les étiquettes \((n,r)\) rendent disjointes les copies des porteurs et fixent
sans ambiguïté chaque origine et destination. La relation
\(\operatorname{Ext}^{\APP,\rm cal}\) transporte la cellule complète et
préserve les deux identités euclidiennes ; la relation
\(\operatorname{Ext}^{\rm fib,cal}\) relie les fibres des états
complets antérieur et postérieur : elle transporte feuillet, orientation, préfixe,
cylindre, frontière et registre et incorpore la retenue, la mémoire et la signature
actualisées. La sélection ne modifie pas le registre et le transport concatène
exactement une entrée et produit la présignature de destination
\eqref{eq:prefirma-transporte-cal-ch}. L’actualisation recalcule la signature au moyen de
l’application typée \(\operatorname{Sig}^{\rm cal}_{\widetilde q}\) à partir du
feuillet, de l’orientation, de la retenue, de la route, de la frontière et de la mémoire déjà déterminés. L’application
\(\operatorname{Upd}^{\rm cal}\) est unaire car l’arête contenue dans
\(z_{\varepsilon,r}(x)\) détermine ses cocycles.

Les identités sont les diagonales et les extensions par mots sont les
compositions de \eqref{eq:ext-palabra-calibrada-ch}. La concaténation des
préfixes et des registres est associative ; la mémoire vérifie la loi de cocycle et
la retenue vérifie
\eqref{eq:ley-carry-correcta-ch}. Par conséquent, les deux parenthésages de trois
arêtes produisent le même état. La formule
\eqref{eq:truncamiento-calibrado-definido-ch} définit la troncature sur
les états complets et l’unicité du dernier registre prouve qu’elle retrouve
l’origine dans tous les champs. Ce sont les lois d’identité, de source,
de but, de composition, de mémoire, de retenue et de troncature du TPK enrichi.
En particulier, les relations d’extension sont déterminées comme graphes
d’applications effectives, et non comme conséquences de simples conditions
nécessaires.
\end{proof}

Chaque \(\gamma_{\varepsilon,k_n}\) comprend exactement neuf applications de \(\operatorname{Upd}^{\rm cal}_{e_{\varepsilon,n,r}}\circ
\mathsf{Tra}^{\rm cal}_{e_{\varepsilon,n,r}}\circ
\mathsf{Sel}^{\rm cal}_{e_{\varepsilon,n,r}}\). La survie sera prouvée au moyen d'une queue construite et ne sera pas utilisée pour définir \(G_n\).

\begin{lemma}[Trisection interne effective d'un tour]
\label{lem:elevacion-catalogal-nonadica}
Pour tout \(n\ge0\), tout \(x\in G_n\) et tout \(\varepsilon\in\{-1,0,+1\}\), l'expression \(\gamma_{\varepsilon,k_n}(x)\) est définie par neuf transitions TPK, avec leurs relations synchronisées \(\operatorname{Ext}^{\APP,\rm cal}\) et \(\operatorname{Ext}^{\rm fib,cal}\), et satisfait
\begin{align}
 \rho^{\rm cal}_{k_n+9,k_n}\gamma_{\varepsilon,k_n}(x)&=x,
 \label{eq:truncamiento-injerto-ch}\\
 g\bigl(\gamma_{\varepsilon,k_n}(x)\bigr)
   &=g(x),&
 \jmath_{k_n+9}\!\left(\operatorname{mem}_{k_n+9}
        (\gamma_{\varepsilon,k_n}(x))\right)
   &=\jmath_{k_n}\!\left(\operatorname{mem}_{k_n}(x)\right)+u,
 \label{eq:fase-memoria-injerto-ch}\\
 c\bigl(\gamma_{\varepsilon,k_n}(x)\bigr)
   &=c(x)+\varepsilon\kappa_{\rm Carr}.&&
 \label{eq:carry-injerto-ch}
\end{align}
\label{eq:tres-vueltas-propiedades-ch}
\label{eq:tres-vueltas-acarreo-ch}
Le registre ordonné des neuf arêtes détermine \(\varepsilon\) de manière unique. Les trois images sont donc disjointes. Chaque extrémité appartient à \(G_{n+1}\) et admet à nouveau les trois applications \(\gamma_{-1,k_{n+1}},\gamma_{0,k_{n+1}},
\gamma_{+1,k_{n+1}}\).
\end{lemma}

\begin{proof}
Les formules \eqref{eq:celula-completa-fase-ch} parcourent une fois les neuf phases et renvoient la copie \(a_{n+1,1}\) de \(a(1)\). À chaque étape, la relation APP transporte le sextuplet complet ; \(\mathsf{Sel}^{\rm cal}\) fixe le type TRIT à partir de la marque de phase \(r\) conservée dans l'état ; et \(\mathsf{Tra}^{\rm cal}\) ajoute exactement un symbole au préfixe, un segment au cylindre et sa paire ordonnée d'extrémités à la frontière. Par récurrence sur \(r\), les tuples explicites \(y_{\varepsilon,r}(x)\) et \(z_{\varepsilon,r}(x)\) satisfont les prédicats des relations indiquées, et la mise à jour unaire produit l'état du niveau suivant. L'identité \(\rho^{\rm cal}_{d_{n,r}+1,d_{n,r}}
(x_{\varepsilon,r}(x))=x_{\varepsilon,r-1}(x)\) est la définition \eqref{eq:truncamiento-calibrado-definido-ch} ; son caractère bien défini découle de l'unicité du dernier registre. La composition des neuf applications de troncature récupère \(x\), ce qui prouve \eqref{eq:truncamiento-injerto-ch} sans inverser séparément les champs écrasés.

La phase \(9\to1\) apparaît une seule fois. Par
\eqref{eq:memoria-arista-ch} et la loi de cocycle, la partie linéaire de mémoire
transporte la mémoire préalable et sa partie translationnelle ajoute exactement \(u\)
dans la limite directe. La compatibilité
\(\jmath_{d+1}\jmath_{d+1,d}=\jmath_d\) prouve
\eqref{eq:fase-memoria-injerto-ch} ; le retour concerne la phase, non
l’état enrichi.

Par \eqref{eq:carry-arista-derivado-ch}, huit arêtes ont un incrément nul et
l’arête qui complète \(\mathfrak p_\varepsilon\) possède l’incrément calculé
\(\chi(\varepsilon,\varepsilon)\kappa_{\rm Carr}\). La loi
\eqref{eq:ley-carry-correcta-ch} et
\eqref{eq:carry-no-parametrico-ch} donnent
\eqref{eq:carry-injerto-ch}. Le registre conserve
\((\mathfrak p_\varepsilon,2\varepsilon,r_3(2\varepsilon),
\chi(\varepsilon,\varepsilon))\) ; les trois valeurs de ce registre sont
distinctes même lorsque l’incrément scalaire est nul. C’est pourquoi les images
ne sont pas identifiées et le décodeur de bloc retrouve un unique
\(\varepsilon\).

Les définitions des arêtes dépendent seulement de la nouvelle phase et concatènent,
au lieu de remplacer, préfixe, cylindre, frontière et registre. Elles sont aussi
invariantes sous
\(\jmath_d\operatorname{mem}_d\mapsto
  \jmath_d\operatorname{mem}_d+nu\)
et sous la translation de la retenue
d’entrée. Par conséquent, la même construction s’applique à l’extrémité de
tout tour. Pour chaque état fini, la suite explicite
\[
 x,\quad\gamma_{0,k_n}(x),\quad
 \gamma_{0,k_{n+1}}\gamma_{0,k_n}(x),\quad\ldots
\]
est une queue infinie compatible. Cela démontre la survie de tous les
états construits et la disponibilité réitérée des trois branches, sans
postuler aucune des deux propriétés. Nous pouvons donc, enfin,
identifier, conformément aux définitions de niveau précédentes,
\begin{equation}
 H_{k_n}^{\rm cal}:=G_n
 \label{eq:subarbol-calibrado-superviviente-ch}
\end{equation}
pour tout \(n\ge0\).
\end{proof}

La clôture précédente distingue les trois horloges impliquées. Neuf émissions
ramènent la phase dans \(C_9\) et font avancer la mémoire. Douze tours, soit
\(108\) émissions, ramènent en outre la position du calendrier observable ;
le registre et la mémoire ne sont pas non plus effacés à ce moment. L’opérateur
\(\mathcal K_{\rm ph}\) agit uniquement ensuite, comme reconnaissance de
mémoire réduite des tours déjà construits ; il n’engendre jamais les arêtes de
\(\operatorname{Ext}^{\rm enr,cal}\).


Soit
\[
 \upsilon:\mathbb F_3\longrightarrow\{-1,0,+1\},\qquad
 \upsilon(0)=0,\quad\upsilon(1)=+1,\quad\upsilon(2)=-1
\]
le représentant équilibré. Pour \(b=(b_1,\ldots,b_6)\in\mathbb F_3^6\), le prolongement de bloc est la composition ordonnée de six tours
\[
 \Gamma_{b,k}^{[54]}:=
 \gamma_{\upsilon(b_6),k+45}\circ\gamma_{\upsilon(b_5),k+36}\circ
 \gamma_{\upsilon(b_4),k+27}\circ\gamma_{\upsilon(b_3),k+18}\circ
 \gamma_{\upsilon(b_2),k+9}\circ\gamma_{\upsilon(b_1),k}.
\tag{C.39d}
\label{eq:gamma-bloque-54-ch}
\]
Au niveau de bloc, on emploie toujours \(k=k_{6N}=1+54N\) ; chaque facteur a donc le domaine construit par le facteur précédent. Fixons un germe APP canonique \(x_\star\in H^{\rm cal}_1\) et définissons récursivement
\[
 \begin{aligned}
  Y_0&:=\{x_\star\},\qquad Y_N\subseteq H_{1+54N}^{\rm cal},\\
  Y_{N+1}&:=\{\Gamma_{b,1+54N}^{[54]}(x):
                 x\in Y_N,\ b\in\mathbb F_3^6\},\\
  \rho^{\rm blk}_{N+1,N}
    &:=\rho^{\rm cal}_{1+54(N+1),1+54N}
       \!\upharpoonright Y_{N+1},\\
  \operatorname{Ext}^{\rm cal}_N(x)
    &:=\{\Gamma_{b,1+54N}^{[54]}(x):b\in\mathbb F_3^6\}.
 \end{aligned}
\tag{C.39e}
\label{eq:sincronizacion-54-ch}
\]
La notation \(\rho^{\rm cal}_{b,a}\) désigne la composition de toutes les suppressions
entre les niveaux \(b\) et \(a\). Par construction, chaque prolongement de
\(\operatorname{Ext}^{\rm cal}_N(x)\) est une chaîne de cinquante-quatre
arêtes TPK, et non un saut du quotient de Markov.

Le lecteur est défini directement sur le registre, avant de faire appel à une représentation au moyen de \(\Gamma_b\). Le lemme précédent fournit un décodeur
\[
 \operatorname{dec}_9:\{\text{segments engendrés de neuf arêtes}\}
 \longrightarrow\{-1,0,+1\},
\]
qui lit le couple TRIT et sa réduction dans le segment. Si
\(\operatorname{Led}^{(9)}_{j,i}(x)\) désigne le \(i\)-ième segment nonadique
du bloc \(j\) du registre de \(x\), on fixe
\begin{equation}
 \beta_{{\rm blk},N}(x):=
 \Bigl(
   \bigl(\upsilon^{-1}\operatorname{dec}_9
          (\operatorname{Led}^{(9)}_{j,i}(x))\bigr)_{i=1}^{6}
 \Bigr)_{j=0}^{N-1}.
 \label{eq:lector-bloques-ch}
\end{equation}
Pour \(N=0\), la valeur est le mot vide. Cette définition ne présuppose pas que la décomposition de \(x\) en composition d'applications \(\Gamma_b\) soit unique.

Spécialisons maintenant la conjugaison arithmétique précédente à
\(p=3\), \(d=6\) et \(A=I_6\). Si
\(\pi_{N+1,N}:(\mathbb Z/3^{N+1}\mathbb Z)^6\to
(\mathbb Z/3^N\mathbb Z)^6\) est la réduction canonique, nous définissons
\begin{equation}
 \Phi_N:=\Psi_N^{-1}\circ\beta_{{\rm blk},N}:
 Y_N\longrightarrow(\mathbb Z/3^N\mathbb Z)^6.
 \label{eq:phi-hensel-bloques-ch}
\end{equation}
Cette définition n’assigne pas une coordonnée après avoir observé une branche : le
lecteur \(\beta_{{\rm blk},N}\) s’obtient à partir du registre TPK déjà construit et
\(\Psi_N^{-1}\) est l’inverse explicite, fixé auparavant, de la tour
arithmétique non filtrée.

\begin{proposition}[Conjugaison finie et recouvrement complet de la fibre]
\label{prop:cobertura-729-seccion}
Pour tout \(N\geq0\),
\[
 \beta_{{\rm blk},N}:Y_N\xrightarrow{\ \cong\ }
 (\mathbb F_3^6)^N
\]
est une bijection compatible avec les troncatures, et pour chaque \(x\in Y_N\)
\[
 \begin{gathered}
 \beta_{{\rm blk},N+1}
 \bigl[\operatorname{Ext}^{\rm cal}_N(x)\bigr]
 =\{\beta_{{\rm blk},N}(x)^\frown b:
       b\in\mathbb F_3^6\},\\
 \operatorname{pref}_N\circ\beta_{{\rm blk},N+1}
 =\beta_{{\rm blk},N}\circ\rho^{\rm blk}_{N+1,N}.
 \end{gathered}
\tag{C.40$_0$}
\label{eq:cobertura-729-ch}
\]
Les applications \(\Phi_N\) sont des bijections et vérifient l’identité
d’entrelacement
\begin{equation}
 \pi_{N+1,N}\circ\Phi_{N+1}
 =\Phi_N\circ\rho^{\rm blk}_{N+1,N}.
 \label{eq:entrelazamiento-hensel-ext-ch}
\end{equation}
Par conséquent, la tour des histoires engendrées et la tour arithmétique de Hensel sont isomorphes en tant que systèmes projectifs finis.
\end{proposition}

\begin{proof}
L’assertion est immédiate pour \(N=0\). Supposons qu’elle soit vraie au niveau
\(N\). Le lemme précédent construit, pour chaque mot
\(b\in\mathbb F_3^6\), un prolongement
    \(\Gamma_{b,1+54N}^{[54]}(x)\). Si \(b\ne b'\), la première coordonnée par laquelle
ils diffèrent détermine une paire TRIT et un cocycle distincts dans le premier
segment nonadique divergent du registre généalogique. Ce registre conserve
les arêtes ordonnées, et non seulement la somme finale de leurs retenues ; les deux
histoires sont donc distinctes même si les retenues de tours ultérieurs se
compensent. La lecture directe \eqref{eq:lector-bloques-ch} donne alors
\begin{equation}
 \beta_{{\rm blk},N+1}
 \bigl(\Gamma_{b,1+54N}^{[54]}(x)\bigr)
 =\beta_{{\rm blk},N}(x)^\frown b.
 \label{eq:beta-concatenacion-demostrada-ch}
\end{equation}
Les six trisections successives produisent exactement \(3^6\) prolongements distincts. La suppression des cinquante-quatre dernières étapes récupère \(x\), et la suppression du dernier bloc du registre récupère \(\beta_{{\rm blk},N}(x)\). Cela prouve simultanément la bijection et la compatibilité de \(\beta_{{\rm blk},N}\) avec les troncatures.

La proposition~\ref{prop:ch-hensel-bruto} prouve, indépendamment, que
\(\Psi_N\) est bijective et que
\[
 \operatorname{pref}_N\circ\Psi_{N+1}
 =\Psi_N\circ\pi_{N+1,N}.
\]
En composant cette égalité avec la compatibilité de \(\beta_{\rm blk}\) et en
appliquant \(\Psi_N^{-1}\), on obtient
\[
 \pi_{N+1,N}\Psi_{N+1}^{-1}\beta_{{\rm blk},N+1}
 =\Psi_N^{-1}\beta_{{\rm blk},N}\rho^{\rm blk}_{N+1,N},
\]
qui est exactement \eqref{eq:entrelazamiento-hensel-ext-ch}. L’égalité ne
provient pas de la définition d’une coordonnée ad hoc sur chaque branche : elle résulte de la composition
de deux systèmes de bijections construits et compatibles à toutes les
profondeurs.
\end{proof}

La limite projective
\[
 \Omega_\Gamma^{\rm cal}:=
 \varprojlim_N(Y_N,\rho^{\rm blk}_{N+1,N})
\tag{C.40a$_0$}
\label{eq:omega-gamma-catalogal-ch}
\]
est la limite d'histoires du modèle TPK construit arête par arête dans le lemme précédent. Sa projection par \(\pi_{\rm obs}\) conserve les configurations du calendrier observable, mais l'affirmation ne dépend pas de l'identification de ce modèle autonome à une relation ambiante définie uniquement par des conditions nécessaires. Les bijections finies passent à la limite et produisent l'homéomorphisme cylindrique
\[
 \Lambda:(\mathbb F_3^6)^{\mathbb N}
 \xrightarrow{\ \cong\ }\Omega_\Gamma^{\rm cal},
 \qquad
 \beta_{\rm blk}:=\Lambda^{-1}.
\tag{C.40b$_0$}
\label{eq:homeomorfismo-calibrado-ch}
\]
La projection canonique est notée
\(\rho^{\rm blk}_{\infty,N}:\Omega_\Gamma^{\rm cal}\to Y_N\).
Ainsi, \(\Lambda\) établit une équivalence compatible avec les
troncatures : chaque ruban ternaire encode une histoire du TPK et chaque
préfixe se prolonge au moyen d’arêtes du même générateur.

% BEGIN REV09_RECTAS27_PROLONGACION
\subsubsection{Prolongement réversible en coordonnées de droites}
\label{corpus9:xii:rectas-prolongacion}

La proposition~\ref{prop:cobertura-729-seccion} a construit les
bijections \(\beta_{{\rm blk},N}\) au moyen de segments nonadiques du
registre ordonné, avec leur compatibilité avec la suppression des
blocs. L’homéomorphisme
\eqref{eq:homeomorfismo-calibrado-ch} identifie la limite de ces
histoires engendrées à la suite de leurs mots de six trits.
Nous composons maintenant ce lecteur avec le système de coordonnées fini
\eqref{corpus9:xii:rectas-codificacion}, en conservant sa loi centrale.

\begin{theorem}[Transport cylindrique des histoires vers des paires de droites]
\label{corpus9:xii:rectas-homeomorfismo}
Pour chaque \(N\geq0\), l’application
\begin{equation}
 \mathcal E_N:=\mathcal E^{\times N}\circ\beta_{{\rm blk},N}:
 Y_N\xrightarrow{\ \cong\ }
 (\mathscr L_{27}\times\mathscr L_{27})^N
 \label{corpus9:xii:rectas-nivel}
\end{equation}
est une bijection. Ces bijections satisfont
\begin{equation}
 \operatorname{pref}_N\circ\mathcal E_{N+1}
 =\mathcal E_N\circ\rho^{\rm blk}_{N+1,N}
 \label{corpus9:xii:rectas-truncamiento}
\end{equation}
et déterminent l’homéomorphisme

\begin{equation}
 \mathcal E_\infty:=\mathcal E^{\mathbb N}\circ\beta_{\rm blk}:
 \Omega_\Gamma^{\rm cal}\xrightarrow{\ \cong\ }
 (\mathscr L_{27}\times\mathscr L_{27})^{\mathbb N}.
 \label{corpus9:xii:rectas-limite}
\end{equation}
Son inverse est
\(\beta_{\rm blk}^{-1}\circ(\mathcal E^{-1})^{\mathbb N}\).

\end{theorem}
\begin{proof}
La proposition~\ref{corpus9:xii:rectas-biyeccion} fournit
l’inverse effectif de chaque facteur \(\mathcal E\), au moyen du tableau
des droites et de l’entrelacement des coordonnées. C’est pourquoi
\(\mathcal E_N^{-1}=\beta_{{\rm blk},N}^{-1}
\circ(\mathcal E^{-1})^{\times N}\). Pour \(N=0\), les deux applications
agissent sur le mot vide. L’application coordonnée par coordonnée
commute avec la suppression du dernier facteur. En utilisant
\eqref{eq:cobertura-729-ch}, on obtient
\[
 \begin{aligned}
 \operatorname{pref}_N\mathcal E^{\times(N+1)}
           \beta_{{\rm blk},N+1}
 &=\mathcal E^{\times N}\operatorname{pref}_N
           \beta_{{\rm blk},N+1}\\
 &=\mathcal E^{\times N}\beta_{{\rm blk},N}
           \rho^{\rm blk}_{N+1,N}.
 \end{aligned}
\]
C’est \eqref{corpus9:xii:rectas-truncamiento}. Une suite de paires
de droites détermine, par l’inverse de chaque facteur, une suite unique
de mots. L’homéomorphisme
\eqref{eq:homeomorfismo-calibrado-ch} reconstruit à partir de celle-ci une
histoire unique dans \(\Omega_\Gamma^{\rm cal}\). Les compositions
récupèrent respectivement l’histoire et la suite initiales.
Dans les topologies cylindriques, un ouvert de base fixe un nombre fini
de coordonnées. L’application et son inverse transforment ces conditions
en conditions sur le même nombre de coordonnées ; toutes deux sont
continues. Cela démontre l’affirmation topologique et la compatibilité
de la limite avec toutes les troncatures finies.

\end{proof}

Pour conserver l’information de l’état complet, nous notons
\(\zeta\) ses registres restants : repère, orientation, retenue,
chemin, mémoire et données de frontière. Sur le domaine admissible déjà
construit, la transformation étendue est
\[
 \widehat{\mathcal E}(w,\zeta)=(\mathcal E(w),\zeta).
\]
Son domaine image conserve les mêmes conditions de compatibilité,
exprimées au moyen de l’inverse explicite de la première coordonnée. Si
\(F:\mathcal X\to\mathcal Y\) est un opérateur entre domaines d’états,
son transport est
\[
 F^{\mathcal E}
 =\widehat{\mathcal E}_{\mathcal Y}\circ F
       \circ\widehat{\mathcal E}_{\mathcal X}^{-1}.
\]
L’identité \(F^{\mathcal E}\widehat{\mathcal E}_{\mathcal X}
=\widehat{\mathcal E}_{\mathcal Y}F\) s’obtient en annulant les deux
applications inverses. Pour un prolongement relationnel
\(\mathcal R\subset\mathcal X\times\mathcal Y\), la définition
correspondante est
\[
 \mathcal R^{\mathcal E}
 =\{(\widehat{\mathcal E}_{\mathcal X}x,
       \widehat{\mathcal E}_{\mathcal Y}y):(x,y)\in\mathcal R\}.
\]
La transformation inverse récupère exactement \(\mathcal R\).
Si le prolongement conserve en outre un témoin \(\omega\) dans un domaine
\(W\), sa réalisation \(\widetilde{\mathcal R}\subset
\mathcal X\times W\times\mathcal Y\) est transportée par
\[
 (x,\omega,y)\longmapsto
 (\widehat{\mathcal E}_{\mathcal X}x,\omega,
       \widehat{\mathcal E}_{\mathcal Y}y).
\]
L’inverse agit sur les deux extrémités et conserve \(\omega\). Ainsi,
des témoins différents restent différents, même s’ils partagent les
extrémités, et la multiplicité enregistrée demeure. La condition de
fonctionnalité se conserve sur les domaines où elle était déjà établie.

La couverture à profondeur arbitraire provient ici du domaine
\(\Omega_\Gamma^{\rm cal}\) et de son lecteur démontré, tandis que les
paires de droites fournissent des coordonnées réversibles de ce domaine.
Le recensement initial de \(243\) mots et la capacité globale de
\(729\) prolongements par bloc conservent leurs fonctions distinctes.
La classification diagonale, incidente et gauchie accompagne chaque facteur ;
les conditions qui sélectionnent des marches adjacentes se transportent comme
restrictions de la suite complète.

% END REV09_RECTAS27_PROLONGACION


Les sous-cylindres sont maintenant définis comme sous-ensembles de la limite, et non comme paires
de suffixes. Pour \(x\in Y_N\) et \(b\in\mathbb F_3^6\), soient
\begin{align}
 \mathscr Z_N(x)
   &:=\{\xi\in\Omega_\Gamma^{\rm cal}:
       \rho^{\rm blk}_{\infty,N}(\xi)=x\},
 \label{eq:cilindro-real-limite-ch}\\
 \mathscr Z_{N+1}(x;b)
   &:=\{\xi\in\Omega_\Gamma^{\rm cal}:
       \rho^{\rm blk}_{\infty,N+1}(\xi)
       =\Gamma_{b,1+54N}^{[54]}(x)\}.
 \label{eq:subcilindro-real-limite-ch}
\end{align}

\begin{theorem}[Partition effective en \(729\) sous-cylindres de raffinement]
\label{thm:particion-729-subcilindros-ch}
Pour tout \(N\ge0\) et tout \(x\in Y_N\), on a l’union disjointe
\begin{equation}
 \mathscr Z_N(x)=
 \bigsqcup_{b\in\mathbb F_3^6}\mathscr Z_{N+1}(x;b).
 \label{eq:particion-real-729-ch}
\end{equation}
Les \(3^6=729\) termes sont des sous-cylindres de raffinement non vides et
ouverts-fermés ; chacun est réalisé par cinquante-quatre arêtes du TPK
construit.
\end{theorem}

\begin{proof}
Soit \(\xi\in\mathscr Z_N(x)\). Sa coordonnée de niveau \(N+1\) appartient à \(Y_{N+1}\), se tronque en \(x\) et, par la bijection \(\beta_{{\rm blk},N+1}\), possède un dernier bloc unique \(b\in\mathbb F_3^6\). L'identité \eqref{eq:beta-concatenacion-demostrada-ch} implique alors \(\rho^{\rm blk}_{\infty,N+1}(\xi)
=\Gamma_{b,1+54N}^{[54]}(x)\) ; par conséquent, \(\xi\) appartient à l'un des termes du membre de droite. L'inclusion inverse est immédiate par troncature. Si deux termes s'intersectaient, leurs coordonnées de niveau \(N+1\) seraient égales et le décodeur du registre donnerait \(b=b'\), de sorte que la réunion est disjointe.

Chaque extrémité \(\Gamma_b^{[54]}(x)\) admet une queue infinie obtenue
en itérant le prolongement de bloc nul ; tous les termes sont donc non
vides. Comme \(Y_{N+1}\) est fini et discret, ce sont des images réciproques de points par
une projection continue de la limite projective et, par conséquent, ils sont ouverts et fermés. La
définition même de \(\Gamma_b^{[54]}\) atteste ses cinquante-quatre arêtes.
\end{proof}

La collection \((\Phi_N)_N\) induit finalement l'homéomorphisme des limites
\begin{equation}
 \Phi_\infty:\Omega_\Gamma^{\rm cal}
 \xrightarrow{\ \cong\ }\mathbb Z_3^6,\qquad
 \Phi_\infty(\xi):=
 \bigl(\Phi_N(\rho^{\rm blk}_{\infty,N}(\xi))\bigr)_{N\ge0},
 \label{eq:phi-infinito-hensel-ext-ch}
\end{equation}
où
\(\mathbb Z_3^6=\varprojlim_N(\mathbb Z/3^N\mathbb Z)^6\).
L’équation \eqref{eq:entrelazamiento-hensel-ext-ch} garantit que le membre
droit est une classe compatible, et les bijections finies prouvent
l’injectivité, la surjectivité et la continuité dans les deux sens. Sous
\(\Phi_\infty\), \(\mathscr Z_N(x)\) est exactement la classe de congruence
\[
 \{z\in\mathbb Z_3^6:z\equiv\Phi_N(x)\pmod{3^N}\},
\]
et les ensembles \(\mathscr Z_{N+1}(x;b)\) sont ses \(729\) classes de raffinement
modulo \(3^{N+1}\). Tel est le raccord effectif entre prolongement TPK,
cylindres de la limite et tour de Hensel.

Pour \(b=(b_1,\ldots,b_6)\in\mathbb F_3^6\), posons
\[
 \nu(b):=\sum_{i=1}^{6}[b_i]_3\,3^{6-i}
 \in\{0,\ldots,728\},
\tag{C.40c$_0$}
\label{eq:valor-bloque-ch}
\]
où \([b_i]_3\in\{0,1,2\}\) est le représentant positionnel. Les
\(243\) mots visibles du recensement initial sont l’image de la première
émission ; ils ne se confondent pas avec la capacité \(3^6\) de chaque bloc de
prolongement. Enfin,
\[
 X_{\infty,\mathbb Z}^{\rm cal}
 :=\mathbb Z_{\rm HMT}\times\Omega_\Gamma^{\rm cal}.
\]
